% GENERATED FILE. Edit canonical lessons, then run npm run generate:books.
% canonical-source-hash: fnv1a64:b8f8fda6eb005d9b
% canonical-lessons: LA-C282-primum-deinde, LA-C282-dum, LA-C282-lectio, LA-R282-a-story-in-order

\chapter{A Story in Order}
\label{ch:la-a-story-in-order}
\hlchaptermodality{282}

\begin{chapteropening}
\textbf{By the end of this chapter:} \emph{I can tell a short story in order, with prīmum, deinde, tum and postrēmō, and hold two things together with dum.}

Complete LA-R282-a-story-in-order: I can tell a short story in order, with prīmum, deinde, tum and postrēmō, and hold two things together with dum.
\end{chapteropening}

\section[prīmum, deinde, tum, postrēmō]{prīmum, deinde, tum, postrēmō — the rungs of a story}

\label{lesson:LA-C282-primum-deinde}



\begin{quote}
Say \textquotedblleft{}I came\textquotedblright{} and \textquotedblleft{}I saw\textquotedblright{}. (\emph{Vēnī. Vīdī.}) The verb for \textquotedblleft{}to win\textquotedblright{} is \emph{vincō}, and its perfect is \textbf{vīcī}, \textquotedblleft{}I won\textquotedblright{}.

Suetonius reports that when Caesar celebrated his victory in Pontus, a placard carried three words: \textbf{vēnī, vīdī, vīcī}. Three events, three perfects, and not one joining word.
\end{quote}

\begin{grammarlens}[title={Grammar Lens: four rungs}]
Caesar could leave the order bare because three short perfects explain themselves. A longer story names its rungs, with four words you already own:

\begin{quote}
\textbf{Prīmum} ad forum ambulāvī. — First I walked to the forum. \textbf{Deinde} amīcum vīdī. — Next I saw a friend. \textbf{Tum} in caupōnam intrāvimus. — Then we went into the inn. \textbf{Postrēmō} domum ambulāvī. — Last of all I walked home.
\end{quote}

Each rung carries a \textbf{perfect}, because each is an event.

\textbf{Deinde} and \textbf{tum} both come out as \textquotedblleft{}then\textquotedblright{}, and they are not twins. \emph{Deinde} is \textquotedblleft{}next\textquotedblright{}, the following step. \emph{Tum} is \textquotedblleft{}then, at that moment\textquotedblright{}: it points at a time, and the scene can stand there with it. \emph{Tum pluēbat}, at that moment it was raining.
\end{grammarlens}

\begin{cousinweb}
\textbf{prīmum} is the neuter of \emph{prīmus}, \textquotedblleft{}first\textquotedblright{}, standing alone as an adverb (Lewis \& Short list it under \emph{primus}): \textquotedblleft{}as the first thing\textquotedblright{}. \emph{Postrēmō} comes from \emph{postrēmus}, \textquotedblleft{}last\textquotedblright{}, too, though it is the ablative, \textquotedblleft{}at the last\textquotedblright{}, and \emph{deinde} is built from \emph{dē} and \emph{inde}, \textquotedblleft{}from there\textquotedblright{}: from that point on.
\end{cousinweb}

\subsection*{Your turn}
Say these aloud:

\begin{itemize}
\raggedright
  \item \emph{vēnī, vīdī, vīcī} — three events, no joiners
  \item \emph{Prīmum ad forum ambulāvī. Deinde amīcum vīdī.}
  \item \emph{Tum in caupōnam intrāvimus. Postrēmō domum ambulāvī.}
  \item a rung with a scene on it — \emph{Tum pluēbat.}
  \item \emph{Amīcī in caupōnam intrāvērunt et vīnum bibērunt.} — then \emph{Amīca fābulam nārrāvit.}
  \item your own morning in four rungs — \emph{prīmum}, \emph{deinde}, \emph{tum}, \emph{postrēmō}, with a perfect on each
\end{itemize}

\begin{tcolorbox}[breakable,colback=teal!4,colframe=teal!35!black,title={Before you move on}]
\textquotedblleft{}First\textquotedblright{}, \textquotedblleft{}next\textquotedblright{}, \textquotedblleft{}then\textquotedblright{}, \textquotedblleft{}lastly\textquotedblright{}? (\emph{\textbf{Prīmum, deinde, tum, postrēmō.}}) Which \textquotedblleft{}then\textquotedblright{} points at a moment and can hold a scene? (\emph{\textbf{Tum.}}) Which tense goes on each rung? (The \textbf{perfect}: each is an event.) What did Caesar's three words say? (\emph{\textbf{Vēnī, vīdī, vīcī}}: I came, I saw, I won.)
\end{tcolorbox}
\section[dum]{dum — while}

\label{lesson:LA-C282-dum}



\begin{quote}
Say \textquotedblleft{}first I walked to the forum\textquotedblright{}. (\emph{Prīmum ad forum ambulāvī.}) Then a scene: \textquotedblleft{}it was raining\textquotedblright{}. (\emph{Pluēbat.})
\end{quote}

\subsection*{You'll want to know: dum}
\textbf{dum} — \textquotedblleft{}while\textquotedblright{}. It opens a stretch of time, and something happens inside it:

\begin{quote}
\textbf{Dum in forō ambulant, amīcam vīdērunt.} — While they were walking in the forum, they saw a friend.
\end{quote}

The seeing is an event, so it is a perfect, \emph{vīdērunt}. The walking is the stretch it happened in, and \emph{dum} holds it open.

\begin{grammarlens}[title={Grammar Lens: dum keeps the present}]
Listen to \emph{ambulant}. The story is in the past, and the verb after \emph{dum} is the \textbf{present}: \textquotedblleft{}they walk\textquotedblright{}. When \emph{dum} means \textquotedblleft{}while\textquotedblright{}, a past story usually keeps the present after it, and English turns it into \textquotedblleft{}were walking\textquotedblright{}.

\begin{quote}
\textbf{Dum mercātōrēs pānem vendunt, amīca vēnit.} — While the merchants were selling bread, my friend came.
\end{quote}

Caesar writes it that way: \emph{dum haec geruntur}, \textquotedblleft{}while these things are being done\textquotedblright{}, and then the event in the perfect (Lewis \& Short quote it from his war in Gaul). In the main sentence, the scene still takes the imperfect: \emph{mercātōrēs pānem vendēbant}.
\end{grammarlens}

\begin{cousinweb}
The origin of \textbf{dum} is not settled. Lewis \& Short connect it with \emph{diū}, \textquotedblleft{}for a long time\textquotedblright{}, the word whose own lesson traced it to the root of \emph{diēs}, \textquotedblleft{}day\textquotedblright{}; EtymDB 2.1 traces \emph{dum} to a Proto-Indo-European particle \emph{*de}. Either way it is a small, old word, and its meaning has stayed \textquotedblleft{}for the stretch of time that\textquotedblright{}.
\end{cousinweb}

\subsection*{Your turn}
Say these aloud:

\begin{itemize}
\raggedright
  \item \emph{dum}
  \item \emph{Dum in forō ambulant, amīcam vīdērunt.}
  \item \emph{Dum mercātōrēs pānem vendunt, amīca vēnit.}
  \item the scene without \emph{dum} — \emph{Mercātōrēs pānem vendēbant.}
  \item one of your own — \emph{Dum} plus a present, then a perfect
\end{itemize}

\begin{tcolorbox}[breakable,colback=teal!4,colframe=teal!35!black,title={Before you move on}]
What does \emph{dum} mean? (\textbf{While.}) In a past story, which tense follows \emph{dum}? (The \textbf{present}: \emph{dum ambulant}.) And the event inside the stretch? (The \textbf{perfect}: \emph{vīdērunt}.) Is the origin of \emph{dum} settled? (\textbf{No.})
\end{tcolorbox}
\section[(reading)]{(reading) — a rainy day, told in both pasts}

\label{lesson:LA-C282-lectio}



\begin{quote}
Your first passage gave you forty-two words in a row, all in the present. This one is twice as long, and it happened yesterday.

Every word in it has come up before this lesson.
\end{quote}

\begin{tcolorbox}[breakable,skin=enhanced,colback=orange!6,colframe=orange!45!black,boxrule=0.5pt,arc=1mm,left=6pt,right=6pt,top=4pt,bottom=4pt,fonttitle=\bfseries,title={Read this}]
\begin{quote}
Herī māne caelum obscūrum erat, et pluēbat. Mārcus in vīllā sedēbat et librum legēbat. Subitō amīcus ad portam vēnit et clāmāvit: \textquotedblleft{}Venī ad forum!\textquotedblright{} Prīmum Mārcus nōn volēbat, quia fessus erat. Deinde tamen vēnit, et ad forum ambulāvērunt.
\end{quote}

\begin{quote}
Dum ambulant, amīcam vīdērunt, et amīca quoque vēnit. In forō mercātōrēs pānem et vīnum vendēbant et multum clāmābant. Tum amīcī in caupōnam intrāvērunt. In caupōnā calor erat. Ibi cēnāvērunt et vīnum bibērunt, et amīca fābulam nārrāvit. Postrēmō domum ambulāvērunt. Iam nōn pluēbat, caelum serēnum erat, et Mārcus laetus erat.
\end{quote}

\emph{Read it:} the passage once without stopping; do not translate as you go

\emph{Read it:} the passage again, and sort every verb into one of two piles: the scene, or what happened
\end{tcolorbox}

\begin{culture}
The events pile is the perfects: \emph{vēnit}, \emph{clāmāvit}, \emph{vēnit}, \emph{ambulāvērunt}, \emph{vīdērunt}, \emph{vēnit}, \emph{intrāvērunt}, \emph{cēnāvērunt}, \emph{bibērunt}, \emph{nārrāvit}, \emph{ambulāvērunt}. Eleven things happened, and \emph{prīmum}, \emph{deinde}, \emph{tum} and \emph{postrēmō} keep them in order.

Everything else is scene: the dark sky, the rain, Marcus reading, the merchants shouting, the warmth of the inn. A story is mostly background.

And \emph{dum ambulant} is a present inside a past story, the stretch of time the meeting happened in.
\end{culture}

\subsection*{Your turn}
\begin{itemize}
\raggedright
  \item \emph{Say it:} the passage aloud, once, without stopping
\end{itemize}

\begin{tcolorbox}[breakable,colback=teal!4,colframe=teal!35!black,title={Before you move on}]
Why is \emph{legēbat} in the imperfect? (The reading \textbf{was already going on} when the friend came: it is scene.)

Why is \emph{ambulant} in the present? (It follows \emph{\textbf{dum}}.)

How did the day end? (\emph{\textbf{Iam nōn pluēbat}}: the rain had stopped, and Marcus was happy.)
\end{tcolorbox}
\section[(review)]{Review — a story, told in order}

\label{lesson:LA-R282-a-story-in-order}



\begin{quote}
Say Caesar's three words. (\emph{Vēnī, vīdī, vīcī.}) Then the four rungs. (\emph{Prīmum, deinde, tum, postrēmō.})
\end{quote}

\begin{grammarlens}[title={What you've built}]
Two chapters, and you can \textbf{nārrāre}, tell a story. The scene in the imperfect, \emph{erat}, \emph{pluēbat}, \emph{sedēbat}. The events in the perfect, one per rung, from \emph{prīmum} to \emph{postrēmō}. And \textbf{dum} with a present for the stretch an event happened in.

\begin{quote}
\emph{Dum in forō ambulō, amīcum vīdī. Prīmum nōn clāmāvī, quia fessus eram.} — While I was walking in the forum, I saw a friend. At first I did not call out, because I was tired.
\end{quote}

In that last sentence the event is negative, \emph{nōn clāmāvī}, and it is still an event; the tiredness around it is scene, \emph{fessus eram}.
\end{grammarlens}

\subsection*{Your turn}
Say these aloud:

\begin{itemize}
\raggedright
  \item the rainy day in four rungs, from memory — \emph{prīmum}, \emph{deinde}, \emph{tum}, \emph{postrēmō}
  \item \emph{Dum in forō ambulō, amīcum vīdī.}
  \item the day's last scene — \emph{Iam nōn pluēbat, et Mārcus laetus erat.}
  \item a day of your own: one line of scene, then three events in order
\end{itemize}

\begin{tcolorbox}[breakable,colback=teal!4,colframe=teal!35!black,title={Before you move on}]
Which rung comes last? (\emph{\textbf{Postrēmō.}}) Which tense follows \emph{dum} in a past story? (The \textbf{present}.) Which tense goes on each rung? (The \textbf{perfect}.) And what did Marcus do first, when his friend called? (\textbf{He did not want to come}: \emph{nōn volēbat}.)
\end{tcolorbox}
